\documentclass[10pt,twocolumn]{article}
\usepackage[letterpaper,top=0.68in,bottom=0.68in,left=0.64in,right=0.64in,columnsep=0.24in]{geometry}
\usepackage{mathptmx}
\usepackage{graphicx}
\usepackage{amsmath,amssymb}
\usepackage{booktabs,tabularx,array,multirow}
\usepackage{microtype}
\usepackage{cite}
\usepackage{url}
\usepackage{fancyhdr}
\usepackage{stfloats}
\usepackage{balance}
\usepackage{float}
\usepackage{titlesec}
\usepackage{caption}
\usepackage{enumitem}
\usepackage[hidelinks]{hyperref}

\setlength{\parindent}{1em}
\setlength{\parskip}{0pt}
\setlength{\textfloatsep}{5pt plus 1pt minus 1pt}
\setlength{\floatsep}{4pt plus 1pt minus 1pt}
\setlength{\intextsep}{4pt plus 1pt minus 1pt}
\setlength{\abovecaptionskip}{2pt}
\setlength{\belowcaptionskip}{0pt}
\setlength{\skip\footins}{4pt}
\renewcommand{\baselinestretch}{1.01}

\titleformat{\section}{\centering\bfseries\scshape\normalsize}{\Roman{section}.}{0.45em}{}
\titleformat{\subsection}{\itshape\normalsize}{\Alph{subsection}.}{0.45em}{}
\titleformat{\subsubsection}{\itshape\normalsize}{}{0em}{}
\titlespacing*{\section}{0pt}{5pt}{2pt}
\titlespacing*{\subsection}{0pt}{4pt}{1pt}
\captionsetup{font=footnotesize,labelfont=bf,labelsep=period,justification=justified,singlelinecheck=false}

\pagestyle{fancy}
\fancyhf{}
\fancyfoot[C]{\scriptsize\thepage}
\renewcommand{\headrulewidth}{0pt}
\renewcommand{\footrulewidth}{0pt}

\newcommand{\vect}[1]{\boldsymbol{#1}}
\newcommand{\figpath}{figures}

\begin{document}

\twocolumn[
\begin{@twocolumnfalse}
\begin{center}
{\fontsize{19}{21}\selectfont SparkSense: Motor Profile Guided Edge Anomaly Detection for AC and DC Motor Condition Monitoring under Operating and Sensor Variability\par}
\vspace{0.12in}
{\large\textbf{K. Rajeswari}$^{*}$\par}
{\small Associate Professor, rajeswari@msec.edu.in\par}
\vspace{0.025in}
{\large\textbf{S. Shrikanth}$^{1}$\par}
{\small Final Year B.E. EEE Student, 311523105043@msec.edu.in\par}
\vspace{0.025in}
{\large\textbf{R. Mageshwar}$^{2}$\par}
{\small Final Year B.E. EEE Student, 311523105023@msec.edu.in\par}
\vspace{0.035in}
{\small Department of Electrical and Electronics Engineering, Meenakshi Sundararajan Engineering College, Chennai, Tamil Nadu 600119, India\par}
{\small $^{*}$Corresponding author. $^{1,2}$Final year undergraduate authors.\par}
\vspace{0.08in}
\end{center}
\end{@twocolumnfalse}
]

\noindent\textbf{Abstract:} SparkSense is an edge monitoring framework for mixed AC and DC motor fleets in facilities where continuous specialist supervision and reliable network access may be unavailable. It selects a motor profile, interprets sensor measurements locally, and stores reviewable alerts during communication outages. This study evaluates the monitoring logic for three phase induction and series DC motors using controlled simulation, then tests the vibration detector on independently collected bearing recordings from the University of Ottawa. A quantized autoencoder generated 1 false alarm among 72 nominal AC cases and 14 among 72 nominal DC cases. It detected every simulated imbalance, but also flagged every healthy case exposed to increased sensor gain or broadband noise. An exploratory comparator that used commanded load, current, shaft speed, and a vibration harmonic ratio detected all injected drag and imbalance cases. Under increased gain, its false alarm count remained equal to its nominal count. On measured vibration recordings, the quantized model and a simple RMS threshold each generated 2 false alarms among 20 healthy recordings and detected 8 of 13 fully faulty recordings at matched load. The results establish a reproducible software baseline, a bounded application scope, and a staged route from the breadboard to a pilot product. The central technical limitation remains that amplitude sensitive anomaly scores cannot separate a mechanical change from an acquisition change without operating and sensor context. The measured replay does not validate the harmonic comparator or a complete SparkSense device.

\vspace{2pt}
\noindent\textbf{Index Terms:} AC motor, DC motor, condition monitoring, edge inference, motor profile, anomaly detection, vibration analysis, quantization.

\section{Introduction}
\noindent{\fontsize{22}{20}\selectfont E}\hspace{-0.08em}\textsc{lectric motors} in small factories often operate without continuous condition monitoring. A local monitor must remain useful when connectivity is intermittent, but it must also distinguish changes in a machine from changes in sensor mounting, gain, load, and speed. This distinction is especially important when one platform is intended for both AC and DC motors because identical signal amplitudes can represent different operating states.

SparkSense was originally proposed as a multimodal node with vibration, electrical, thermal, and optional acoustic inputs, a microcontroller for local inference, local storage, and remote reporting. The present study evaluates only the functions supported by reproducible simulation, host tests, and independently collected vibration recordings. It asks a focused and testable question: can motor profile guided inference remain interpretable when motor type, mechanical demand, or measurement conditions change?

Motor fault diagnosis studies can report optimistic performance when correlated records from the same bearing or operating run enter both training and testing partitions \cite{wheat2024}. Integer quantization reduces model storage and can enable efficient inference, but it does not establish fault causality \cite{jacob2018}. Recent TinyML motor monitors demonstrate that embedded vibration inference is feasible in controlled tests \cite{arciniegas2025}, while broader reviews emphasize the continuing importance of domain shift, operating condition, and independent validation \cite{lei2020}.

This paper makes five contributions. First, it defines separate induction and series DC profiles and exposes the assumptions used to generate their electrical and mechanical states. Second, it compares a simple RMS threshold, floating and integer autoencoders, and a load conditioned harmonic comparator under matched disturbances. Third, it uses bearing identity based partitions for independently collected vibration recordings. Fourth, it tests the local event queue during a simulated outage. Fifth, it translates the evidence into a four state operational policy, an explicit application boundary, a target user profile, a commissioning plan, and a budgeted product development route. The analysis distinguishes software evidence from physical device evidence and reports counts with their denominators instead of a single pooled accuracy.

The investigation is organized around four research questions. RQ1 asks whether a detector trained on healthy observations preserves specificity when the same motor operates at different commanded loads. RQ2 asks whether profile specific routing changes the balance between healthy false alarms and simulated fault sensitivity. RQ3 asks whether integer quantization preserves decisions produced by the floating model. RQ4 asks whether conclusions drawn from the controlled generator remain visible when the vibration path is replayed with recordings obtained from an independent test rig. These questions are narrower than complete predictive maintenance because no remaining useful life is estimated and no maintenance intervention is predicted.

The original contribution is the evaluation structure rather than a new neural network. Motor profile, operating command, acquisition disturbance, and asset identity are treated as separate sources of variation. This structure exposes cases in which a detector is sensitive but not specific. It also prevents simulation results, measured replay, and proposed hardware functions from being presented as the same form of evidence.

\section{Related Work}
\subsection{Machine condition monitoring}
Vibration, motor current, temperature, and acoustic measurements reveal different aspects of rotating machine condition. Classical methods use time statistics, spectral energy, shaft order components, and envelope features. Learning methods can combine these features, but their reported performance depends on asset identity, sensor position, load, speed, and data partitioning \cite{lei2020,wheat2024}. ISO 20816 provides general guidance for evaluating machine vibration, but it does not convert a generic anomaly score into a fault diagnosis \cite{iso20816}.

\subsection{Edge inference and TinyML}
TinyML moves preprocessing and inference close to a sensor, which can reduce communication demand and preserve local operation during outages \cite{warden2019}. Quantized networks can reduce weight storage by representing weights and activations as integers \cite{jacob2018}. A valid device claim still requires measured memory, timing, power, sampling behavior, and numerical agreement on the target hardware. This study evaluates host based integer arithmetic and therefore reports algorithmic agreement only.

\subsection{Position of the present work}
Published edge monitors commonly emphasize laboratory classification accuracy, response time, or system integration \cite{arciniegas2025}. SparkSense instead tests a failure mode that an aggregate accuracy can hide: a healthy gain or noise shift can trigger the same decision as a simulated mechanical fault. Table~\ref{tab:position} states the evidence boundary explicitly.

Signal based reviews of induction motor diagnosis group current, vibration, acoustic, thermal, and flux measurements according to the fault mechanisms they can reveal \cite{gangsar2020}. Deep learning reviews likewise emphasize that model capacity does not remove the need for representative operating conditions and independent evaluation \cite{zhao2019}. TensorFlow Lite Micro provides a practical inference runtime for microcontrollers, but its published resource results cannot be transferred to a different network and board without measurement \cite{david2021}. These observations motivate the conservative comparison used here: every learned detector is assessed beside an amplitude threshold and every claimed hardware function is separated from the host computation that was actually executed.

\begin{table}[H]
\caption{Evidence available for each SparkSense function}
\label{tab:position}
\centering
\scriptsize
\begin{tabularx}{\columnwidth}{@{}p{0.22\columnwidth}X p{0.28\columnwidth}@{}}
\toprule
Function & Role & Evidence in this study \\
\midrule
Vibration & Amplitude, band power, and shaft order features & Simulated signals and Ottawa replay \\
Current and speed & Electrical and mechanical context selected by motor profile & Simulated AC and DC states \\
Temperature & Slow contextual trend & Steady state proxy only \\
Event reporting & Local queue followed by remote forwarding & Deterministic host test \\
Physical node & Sensors, microcontroller, storage, and communication & Assembly photograph only \\
\bottomrule
\end{tabularx}
\end{table}

\begin{table*}[t]
\caption{Position of SparkSense relative to representative condition monitoring and edge inference literature}
\label{tab:literature}
\centering
\footnotesize
\begin{tabularx}{0.96\textwidth}{@{}p{0.17\textwidth}p{0.16\textwidth}p{0.17\textwidth}p{0.18\textwidth}X@{}}
\toprule
Reference & Main signal or platform & Evidence source & Principal emphasis & Limitation relevant to this study \\
\midrule
Lei \emph{et al.} \cite{lei2020} & Multiple machinery signals & Review of published studies & Machine learning diagnosis roadmap & Generalization depends on data and operating domain \\
Gangsar and Tiwari \cite{gangsar2020} & Current, vibration, flux, acoustic, thermal & Review of induction motor studies & Signal based motor fault diagnosis & Sensor and fault mechanisms differ across modalities \\
Zhao \emph{et al.} \cite{zhao2019} & Machinery health data & Review and comparative discussion & Deep learning for machine health & Large models do not replace representative validation \\
David \emph{et al.} \cite{david2021} & Microcontroller inference & Embedded runtime evaluation & TensorFlow Lite Micro implementation & Device resource results are platform and model dependent \\
Arciniegas \emph{et al.} \cite{arciniegas2025} & Motor vibration and TinyML & Laboratory IoT case study & Embedded anomaly detection and reporting & Reported performance is tied to the tested laboratory configuration \\
Present study & Motor profile, current, speed, and vibration & Controlled generator plus independent vibration replay & Disturbance specificity and evidence separation & No synchronized physical AC and DC motor validation \\
\bottomrule
\end{tabularx}
\end{table*}

\section{SparkSense Architecture}
SparkSense assigns each asset a nameplate profile and makes the first alarm decision locally. Figure~\ref{fig:architecture} separates the proposed physical interfaces from the tested software paths. The original STM32F411 node, accelerometer, Hall effect current channel, AC energy meter, PT100 sensor, optional microphone, local storage, MQTT transport, database, and dashboard remain design elements. AC energy metering is not treated as a DC measurement. A shaft speed estimate or tachometer is an explicit input when harmonic features are used.

\begin{figure}[t]
\centering
\includegraphics[width=\columnwidth]{\figpath/rId22.png}
\caption{Functional architecture of SparkSense. The motor profile selects the applicable electrical and mechanical context before local feature extraction and anomaly scoring. The event queue retains decisions during a communication outage and forwards them after reconnection. Solid functional blocks represent the proposed system; quantitative evaluation in this study covers the simulated signal path and the host based queue policy.}
\label{fig:architecture}
\end{figure}

The laboratory breadboard in Fig.~\ref{fig:prototype} records an early stage of hardware integration. It shows the sensing and controller assembly used during development, but the available video does not provide synchronized reference measurements or labelled fault trials. The photograph is therefore retained as design documentation, while the quantitative analysis uses the controlled generator and the Ottawa replay.

\begin{figure}[t]
\centering
\includegraphics[width=0.92\columnwidth]{\figpath/image1.png}
\caption{SparkSense laboratory prototype during early integration. The breadboard combines the controller, interface wiring, and indicator hardware shown in the development video. Quantitative results reported in this paper are obtained from the stated simulations and independent vibration recordings rather than from this photograph.}
\label{fig:prototype}
\end{figure}

The original 200 samples per second setting implies a Nyquist frequency of 100 Hz and cannot support features above that frequency. The computational waveform channel therefore uses 2000 samples per second and reports vibration content from 100 to 799 Hz. Physical deployment must choose the accelerometer output rate and filter together \cite{adxl355}. Electrical metering, thermal polling, and waveform acquisition require separate rates. A one second output interval is not a sensor sampling rate.

\section{Application Scope and Intended Users}
The initial SparkSense idea was a common monitoring node for facilities that operate several small and medium motors but cannot install a proprietary monitoring system on every asset. Its intended decision is a condition alert that directs a technician to inspect a machine and its sensor. It is not an automatic trip relay, a certified protection device, or a remaining life estimator. This distinction defines both the target audience and the evidence required before deployment.

The primary users are maintenance teams in small and medium manufacturing plants, motor service centres, campus utility groups, water and wastewater operators, and equipment integrators. Candidate assets include pumps, fans, compressors, conveyors, machine tools, and laboratory or service shop DC drives. The proposed benefit is one configurable workflow for mixed motor types, local operation during network interruption, and event records that preserve the operating context behind an alert. Table~\ref{tab:applications} converts this broad idea into bounded use cases.

\begin{table*}[t]
\caption{Intended SparkSense applications and decision boundaries}
\label{tab:applications}
\centering
\footnotesize
\begin{tabularx}{0.96\textwidth}{@{}p{0.21\textwidth}p{0.25\textwidth}p{0.22\textwidth}X@{}}
\toprule
User, asset, and problem & Required measurements & SparkSense output & Boundary of use \\
\midrule
Plant maintenance; intermittent inspection of pumps, fans, and conveyors in mixed motor fleets & Vibration, current, speed or command, temperature, and sensor status & Profile specific condition alert with evidence record & Technician review is required before maintenance or shutdown \\
Utility or campus operator; distributed ventilation and water assets with weak network coverage & Vibration, current, operating state, timestamp, and local storage status & Local alert retention and later forwarding & No safety protection, process control, or certified availability claim \\
Motor service centre; repeatable post repair baseline for AC and DC machines & Nameplate profile, controlled load cells, current, speed, and vibration & Commissioning baseline and comparison report & Valid only for measured mounting and operating cells \\
Equipment integrator; configurable monitoring module for an original equipment manufacturer pilot & Declared sensor interfaces, motor profile, firmware version, and communication status & Versioned event stream for dashboard or maintenance system & Product integration needs electromagnetic, electrical, and environmental qualification \\
\bottomrule
\end{tabularx}
\end{table*}

The target asset should have repeatable stable operating states, accessible nonhazardous sensor mounting, and a maintenance process able to investigate alerts. Variable speed assets require measured speed or a trustworthy controller command. SparkSense is unsuitable as the sole monitor for personnel safety, hazardous area compliance, emergency shutdown, or high consequence protection. Those applications require certified instrumentation, redundancy, and an engineering safety assessment beyond this study.

\section{Methods}
\subsection{Motor profiles}
The generator represents steady operating points and omits startup, switching transients, saturation, and inverter harmonics. For the balanced three phase induction profile, the simplified per phase equivalent neglects magnetizing current. With RMS voltage $V$, stator resistance $r$, referred rotor resistance $R$, leakage reactance $X$, synchronous angular speed $\omega_s$, and slip $s$, the load torque and rotor speed are
\begin{equation}
T_L=\frac{3V^2(R/s)}{\omega_s\left[(r+R/s)^2+X^2\right]},\qquad
\omega=(1-s)\omega_s .
\label{eq:ac}
\end{equation}
The low slip branch is used and the phase current magnitude is $V/\sqrt{(r+R/s)^2+X^2}$.

For the series DC profile, unsaturated flux is taken as proportional to current and losses other than series resistance are neglected:
\begin{equation}
T_L=KI^2,\quad V=R_sI+KI\omega,\quad
I=\sqrt{T_L/K},\quad \omega=\frac{V-R_sI}{KI}.
\label{eq:dc}
\end{equation}
Base loads are 0.55, 1.00, and 1.45 N m. The AC profile uses 45 V per phase, 50 Hz, a 25 Hz synchronous shaft frequency, and nominal parameters $r=0.72~\Omega$, $R=0.50~\Omega$, and $X=2.3~\Omega$. The DC profile uses 24 V, $R_s=1.2~\Omega$, and $K=0.10$ in the adopted SI equations. Three independent factors drawn uniformly from 0.88 to 1.12 perturb circuit and vibration parameters for each virtual unit. These are illustrative machines, not calibrated replicas.

\subsection{Observations and interventions}
Each case contains one second of current and vibration. AC current is sinusoidal. DC current contains 3.5 percent ripple at six times shaft frequency. At shaft angle $\theta=\omega t$, nominal vibration is
\begin{equation}
x(t)=\left(0.14c+0.02T_L\right)\sin(\theta+0.1)
+0.045\sin(2\theta)+\epsilon(t),
\label{eq:vib}
\end{equation}
where $c$ is a unit specific vibration factor and $\epsilon$ is Gaussian noise with standard deviation 0.026 in arbitrary amplitude units. Drag increases effective torque by 30 percent without changing the nominal command. Imbalance adds a first shaft order component with amplitude 0.17. Healthy stress tests apply a vibration gain of 1.65 or additional broadband noise with standard deviation 0.125. Paired base draws are reused across interventions. These magnitudes are test conditions, not measured fault severities or occurrence probabilities.

\subsection{Local and remote inference}
The local input is current RMS, shaft frequency, a temperature proxy, vibration RMS, and the logarithm of the vibration power fraction from 100 to 799 Hz. Vibration is demeaned, a Hann window is applied to limit spectral leakage \cite{harris1978}, and the five features are standardized using healthy fitting data only. The autoencoder has five inputs, two rectified hidden units, and five linear outputs. For standardized input $\vect{z}$ and reconstruction $\hat{\vect{z}}$, an alarm occurs when
\begin{equation}
S(\vect{z})=\frac{1}{5}\lVert\vect{z}-\hat{\vect{z}}\rVert_2^2>\tau .
\label{eq:score}
\end{equation}
The threshold is the larger empirical 95th percentile from separate healthy calibration groups. Shared and profile specific models are compared with a simple upper vibration RMS threshold. In the host quantizer, weights and input and hidden activations use signed eight bit values, while matrix products accumulate in 32 bit integers. Biases use accumulator units. Scale application and final scoring remain floating point. TensorFlow Lite Micro is a suitable deployment runtime \cite{david2021}, but no claim is made here for its timing, SRAM, power, or numerical equivalence on the proposed board.

A remote Random Forest with 80 trees is trained separately for each motor profile on labelled synthetic healthy, drag, and imbalance cases \cite{breiman2001}. It is evaluated on every nominal test example. Its confusion matrix therefore measures standalone generator classification, not complete alert recall. The AI4I dataset is not used because its tabular labels do not match the waveform inputs or motor states evaluated here.

\subsection{Operational decision policy}
A deployable monitor needs more than an anomaly threshold because invalid acquisition and operation outside the commissioned domain can produce large scores. SparkSense therefore specifies four mutually exclusive output states: valid normal, valid anomaly, sensor or acquisition issue, and outside commissioned operation. Let $q_t$ be one only when sampling continuity, clipping, sensor range, and required context checks pass for window $t$. Let $c_t$ be the current load and speed cell, and let $\mathcal{C}_p$ be the commissioned cells for motor profile $p$. The valid anomaly indicator is
\begin{equation}
a_t=\mathbb{1}\!\left[q_t=1,\ c_t\in\mathcal{C}_p,\ S(\vect{z}_t)>\tau_{p,c_t}\right].
\label{eq:policy}
\end{equation}
If $q_t=0$, the event is labelled as an acquisition issue rather than a motor fault. If $c_t\notin\mathcal{C}_p$, SparkSense abstains and requests additional healthy commissioning data. Only a valid window inside the commissioned domain can be labelled normal or anomalous.

An operational alert may require $m$ anomalous decisions among the most recent $n$ valid windows:
\begin{equation}
E_t=\mathbb{1}\!\left[\sum_{j=0}^{n-1}a_{t-j}\ge m\right].
\label{eq:persistence}
\end{equation}
The pair $(m,n)$ controls persistence and latency and must be fixed using time ordered pilot data. It is not selected from the balanced cases reported here. Every event record retains $q_t$, $c_t$, the score, threshold, model version, and reason code so that an alert can be audited. This policy is a software specification derived from the observed disturbance failures; its field effectiveness remains to be measured.

\subsection{Partitions and uncertainty}
For each profile, 12 virtual units supply fitting data, 8 different units supply healthy calibration, and 12 further units supply test cases. Three loads and two repeats produce 72 test cases for each state and profile. Unit parameters remain fixed across records. Counts are descriptive because 72 records represent repeated measurements from 12 virtual units, not 72 independent machines. Block bootstrap intervals resample the 12 test units and hold fitted weights fixed.

\subsection{Healthy commissioning with a limited budget}
A practical installation cannot assume that every future load and speed has been observed during commissioning. The commissioning record must therefore describe operating coverage, not merely contain many adjacent windows from one run. For each motor, SparkSense stores the motor profile, rated quantities, sensor location, sensor gain, acquisition rate, shaft speed source, and the commanded operating state. Healthy windows are then grouped by stable operating cell. In the present generator the cells are the three commanded loads. In a variable speed installation the cell definition must include speed as well as load.

The limited budget rule used for evaluation is asset separated and session aware. Windows from a single virtual unit remain together, and fitting, calibration, and testing units never overlap. Two repeated observations at each of three loads give six observations per virtual unit. Eight calibration units therefore supply 48 healthy calibration observations per profile. This is not claimed as a universal optimum. It is the smallest evaluated budget in this study that covers every modelled load with repeated observations while preserving separate units for fitting and testing. Collecting more adjacent windows from one unchanged run would increase the record count without increasing coverage of mounting, restart, or operating variability.

For field commissioning, the same principle leads to a staged procedure. First, record healthy operation at the lowest safe stable load, the normal production load, and the highest routinely used load. Second, repeat these measurements after at least one shutdown and restart so that the calibration includes a new acquisition session. Third, inspect current, speed, vibration RMS, clipping, and missing samples before fitting a threshold. Finally, reserve a later healthy session for acceptance testing. If only one operating state is available, the resulting detector must be labelled valid for that state rather than presented as a fleet wide model.

\subsection{Evaluation criteria}
The primary results are recording or case counts with explicit denominators. For nominal healthy cases, specificity is one minus the false alarm fraction. For imposed drag and imbalance, sensitivity is the alarm fraction. Quantization agreement is the fraction of decisions unchanged relative to floating inference. The cloud classifier is summarized with class counts because all classes are balanced by construction. A single pooled accuracy is not used for the local detector because it can conceal both healthy false alarms and missed faults.

The disturbance tests are treated as negative controls. Increased gain changes amplitude without changing the simulated mechanical state, while added broadband noise changes the measurement quality. A detector that alarms on these cases may be sensitive to change, but the alarm cannot be attributed to a mechanical fault. This distinction is central to the interpretation of Table~\ref{tab:simulation} and Fig.~\ref{fig:synthetic}.

\subsection{Exploratory harmonic comparator}
At known shaft speed, least squares sine and cosine fits estimate first and second shaft order amplitudes, $A_1$ and $A_2$. For common positive sensor gain $g$,
\begin{equation}
\frac{gA_1}{gA_2}=\frac{A_1}{A_2},\qquad A_2\ne0.
\label{eq:ratio}
\end{equation}
The ratio can preserve a relative first order change while cancelling common gain. It fails under clipping, frequency dependent response, or unreliable speed. The comparator uses the largest standardized deviation among current RMS, shaft frequency, and the logarithm of $A_1/A_2$. Healthy statistics and thresholds are fitted separately by motor profile and commanded load. The rule was designed after inspection and is structurally matched to the generator. It is an explanatory comparator, not an independently validated fault detector.

\subsection{Measured vibration replay}
The University of Ottawa archive provides vibration recordings from identified bearings under constant load and speed \cite{sehri2023data,sehri2023dataset}. The available archive contained 56 files from 20 bearings. Acceleration sampled at 42 kHz was filtered and resampled to 2 kHz, then divided into ten one second windows. Median window RMS and logarithmic band power fraction produced one vector per recording. Each fold excluded every state of its test bearing, fitted on ten other healthy bearings, and set its threshold using the maximum score from nine further healthy bearings. A two input, one hidden unit model was used because simultaneous SparkSense current, speed, and temperature were unavailable. Four cage fault files for bearings 19 and 20 were absent. Ball fault records were evaluated separately because their applied load differed from the matched load task.

\section{Results}
\subsection{Simulated local decisions}
Table~\ref{tab:simulation} reports alarms among 72 cases in every cell. Healthy, gain, and noise alarms are false positives. Drag and imbalance alarms are detections. The shared quantized model produced no nominal AC false alarm but detected only 13 drag cases. Profile routing increased drag sensitivity to 23 cases with one AC false alarm. In DC, profile routing produced 14 nominal false alarms and 18 drag detections.

\begin{table*}[t]
\caption{Local alarms among 72 simulated cases per condition}
\label{tab:simulation}
\centering
\footnotesize
\begin{tabular}{lccccc}
\toprule
Profile and detector & Healthy & Drag & Imbalance & Healthy gain change & Healthy noise \\
\midrule
AC profile RMS threshold & 2 & 6 & 72 & 72 & 72 \\
AC shared quantized autoencoder & 0 & 13 & 72 & 72 & 72 \\
AC profile quantized autoencoder & 1 & 23 & 72 & 72 & 72 \\
AC exploratory harmonic comparator & 3 & 72 & 72 & 3 & 22 \\
DC profile RMS threshold & 7 & 20 & 72 & 72 & 72 \\
DC shared quantized autoencoder & 7 & 14 & 72 & 72 & 72 \\
DC profile quantized autoencoder & 14 & 18 & 72 & 72 & 72 \\
DC exploratory harmonic comparator & 12 & 72 & 72 & 12 & 31 \\
\bottomrule
\end{tabular}
\end{table*}

Quantization changed none of 360 AC decisions and 6 of 360 DC decisions relative to floating inference. It did not resolve disturbance sensitivity. Resampling test units as blocks with 10,000 draws gave conditional 95 percent intervals of 0 to 4.2 percent for nominal AC false alarms and 6.9 to 31.9 percent for nominal DC false alarms. Drag detection intervals were 29.2 to 33.3 percent for AC and 15.3 to 37.5 percent for DC. These intervals omit generator model error and field uncertainty.

The harmonic comparator detected all drag and imbalance cases. Its gain stress counts equalled its nominal healthy counts, consistent with Eq.~\eqref{eq:ratio}; broadband noise still caused 22 AC and 31 DC alarms. A notification therefore remains an unexplained anomaly until sensor condition and operating context are checked.

The commissioning result is asymmetric across profiles. The quantized profile model produced one nominal false alarm for AC but 14 for DC, even though both profiles used the same number of fitting, calibration, and test units. One DC virtual unit contributed four of its six nominal alarms. This concentration shows why windows must be grouped by asset or session during evaluation. A random window split could distribute observations from that unit across fitting and testing and make the detector appear more stable than it is on a previously unseen unit.

The results also define what additional healthy data would be valuable. Repeating the same nominal run would mainly reduce uncertainty about within run noise. A more informative measurement budget would add new restarts, mounting checks, and operating cells at the limits of normal load and speed. These observations would test whether a threshold represents the motor or only the initial acquisition session. The present data do not establish the exact number of sessions needed for a field deployment, so the proposed commissioning procedure should be treated as a minimum test design rather than a certified requirement.

\begin{figure*}[t]
\centering
\includegraphics[width=0.91\textwidth]{\figpath/rId27.png}
\caption{Alarm counts for matched simulated operating cases. Each cell contains the number of alarms among 72 observations. Columns labelled Healthy, Gain, and Noise quantify false alarms, while Drag and Imbalance quantify detection under the imposed interventions. The figure shows why disturbance tests must be reported together with fault sensitivity.}
\label{fig:synthetic}
\end{figure*}

The remote classifier assigned all AC examples to their generated classes. For DC, four drag examples were labelled healthy and the other 68 drag examples, all healthy examples, and all imbalance examples received their generated labels. This is a consistency result within one generator family. It is not an industrial classification rate.

\subsection{Measured Ottawa replay}
Figure~\ref{fig:ottawa} and Table~\ref{tab:ottawa} report recording level decisions. For matched load fully faulty bearings, the quantized detector produced 18 true negatives, 2 false positives, 8 true positives, and 5 false negatives. The RMS threshold produced the same healthy false alarm and fully faulty detection counts. A pooled accuracy would hide the five missed fully faulty recordings.

\begin{figure*}[t]
\centering
\includegraphics[width=0.88\textwidth]{\figpath/rId30.png}
\caption{Recording level decisions obtained from the University of Ottawa vibration replay. The healthy group contains 20 recordings and the developing and fully faulty groups contain 13 recordings each at matched load. Healthy percentages are false alarm rates. Fault percentages are detection rates. The unloaded ball fault group is evaluated separately and is not included in these bars.}
\label{fig:ottawa}
\end{figure*}

\begin{table}[t]
\caption{Measured recording level decisions}
\label{tab:ottawa}
\centering
\scriptsize
\begin{tabularx}{\columnwidth}{@{}Xccc@{}}
\toprule
Detector & Healthy false alarms & Developing faults & Fully faulty \\
\midrule
Vibration RMS & 2/20 & 8/13 & 8/13 \\
Floating autoencoder & 4/20 & 6/13 & 8/13 \\
Quantized autoencoder & 2/20 & 7/13 & 8/13 \\
\bottomrule
\end{tabularx}
\end{table}

Healthy recordings are reused in the fitting or calibration set of other folds. The results therefore do not provide an independent prospective estimate of field specificity. The measured data also lack simultaneous motor current, command, sensor gain labels, and comparable DC states, so they cannot validate the harmonic comparator.

\subsection{Offline event queue}
A separate test generated one event per second for 60 seconds and imposed a 20 second outage. Incoming events were discarded only when the first in first out queue was full. On reconnection, the queue drained immediately in timestamp order. Eight slots delivered 48 of 60 events, while 32 slots delivered all 60. This verifies the specified queue policy under instantaneous drain. It does not test a live broker, TLS, nonvolatile memory endurance, throughput, end to end latency, or uptime.

The queue test also clarifies the information that must accompany an alarm. Each stored event requires a timestamp, motor identifier, motor profile, operating cell, feature vector, score, threshold, detector version, sensor status, and reason code. Storing only a binary alarm would prevent later review of whether the event was associated with a load transition, a gain change, missing samples, or a persistent condition. A production implementation should use a monotonic sequence number in addition to wall clock time so that ordering remains recoverable after clock correction.

\section{Discussion}
The main result is not a high aggregate accuracy. It is a clear identifiability limit. An amplitude sensitive detector can react in the same way to a simulated mechanical change and to an acquisition disturbance. Within the controlled generator, a shaft order ratio cancelled common gain only when speed was known and the second harmonic remained informative. A practical monitor must therefore record mounting, gain, speed source, command validity, and sensor health. When those inputs are missing, it should report an unexplained anomaly instead of assigning a fault cause.

The measured replay provides external corroboration for the evaluation procedure, but it also weakens the case for unnecessary model complexity. On the fully faulty matched load task, the quantized autoencoder did not outperform the RMS threshold. The autoencoder may become useful when genuinely multimodal measurements are available, but that benefit has not been demonstrated by the Ottawa replay.

This finding changes the role of the learned model in SparkSense. The model should be introduced only when it improves a prespecified decision over a transparent baseline under the same partitions. Complexity is justified by better robustness, earlier detection, or useful discrimination, not by the presence of a neural network alone. Until synchronized multimodal records are available, vibration RMS remains a necessary reference method and the autoencoder remains an experimental component.

The harmonic ratio illustrates the value and the danger of physics informed features. Its common gain cancellation follows directly from Eq.~\eqref{eq:ratio}, so its behavior under the simulated gain change is interpretable. At the same time, the ratio relies on a measurable second shaft order, stable speed, and approximately common gain at both frequencies. These assumptions can fail when the sensor clips, the mounting transfer function changes, or the second order component approaches the noise floor. A field implementation must test these validity conditions before using the ratio in a decision.

SparkSense remains relevant to facilities with mixed motor types and intermittent connectivity, but no savings, industrial equivalence, or operational reliability follows from a block diagram or host simulation. The development budget in Section~\ref{sec:product} is an engineering planning range rather than a measured installed cost or commercial quotation. A fair comparison with commercial or research systems requires the same assets, labels, sensor positions, operating states, and maintenance decisions. Target hardware claims require firmware, a bill of materials, measured memory, timing, power, and sampling verification.

The contribution is therefore an auditable monitoring method and evaluation protocol. It combines motor profile routing, disturbance controls, unit separated partitions, quantization agreement, independent vibration replay, and outage aware event records in one reproducible study. None of these elements alone is new. Their integration addresses a practical weakness in low cost monitoring studies: software behavior, device capability, and field effectiveness are often discussed together even when only one has been measured. The paper keeps those claims separate and identifies the measurements required to move from one evidence level to the next.

The next validation stage should use synchronized current, shaft speed, command or load, vibration, and sensor status from independently identified healthy and faulty AC and DC motors. At least one complete healthy session must remain unused until the analysis is fixed. Fault trials must record how the condition was introduced and verified. Hardware evaluation must report the compiled model size, peak memory, inference time distribution, acquisition timing, energy consumption, queue recovery rate, and numerical agreement with the host implementation. These measurements would resolve the remaining validation questions without relying on inferred specifications.

The study is limited by shared generator assumptions across virtual units, approximate electrical and thermal models, retrospective vibration replay, unavailable cage recordings, no independent DC experiment, and absent single phase testing. It omits startup, saturation, inverter spectra, mounting resonance, electromagnetic compatibility, enclosure effects, and sensor ageing. The results therefore establish a reproducible software baseline and a measurement plan, not a deployable product validation.

\section{Prototype to Product Plan}
\label{sec:product}
Product development should preserve the evidence boundary used in this paper. A larger enclosure, a dashboard, or a production printed circuit board does not compensate for unstable specificity. Each stage in Table~\ref{tab:roadmap} therefore has a technical exit criterion. The cost ranges are author planning estimates in Indian rupees for an academic and small enterprise development programme. They include equipment access, fabrication, and a contingency allowance, but exclude salaries, taxes, building work, and volume manufacturing. Supplier quotations are required before procurement.

\begin{table*}[t]
\caption{Budgeted route from the present breadboard to a pilot ready product}
\label{tab:roadmap}
\centering
\footnotesize
\begin{tabularx}{0.96\textwidth}{@{}p{0.16\textwidth}p{0.28\textwidth}p{0.14\textwidth}X@{}}
\toprule
Stage and deliverable & Main work & Planning budget & Required evidence before exit \\
\midrule
Engineering alpha: reproducible laboratory node & Revised sensor interfaces, protected power input, printed circuit board, enclosure, firmware acquisition, local storage, and three nodes & INR 0.8 to 1.4 lakh & Sampling and filter verification, numerical agreement with the host pipeline, memory, timing, power, data loss, and sensor fault tests \\
Controlled validation rig: engineering validation report & Instrumented induction and DC motors, guarded load system, speed and torque references, calibrated acquisition, mounting fixtures, and repeatable healthy and seeded fault sessions & INR 2.5 to 4.5 lakh & Locked protocol on independent machines and sessions; synchronized current, speed, load, vibration, and sensor state; stable specificity beside simple baselines \\
Limited field pilot: pilot ready system & Eight to ten rugged nodes, installation materials, gateway, database, dashboard, spares, and supervised operation on noncritical assets & INR 1.8 to 3.2 lakh & Alert review log, missed data and outage recovery, drift, maintenance actions, false alerts per operating hour, and user acceptance findings \\
Preproduction: manufacturing package & Design for manufacture, 25 unit build, electromagnetic and electrical precompliance, environmental tests, documentation, and design correction & INR 4.0 to 8.0 lakh & Repeatable production tests, traceable calibration, enclosure and connector verification, cybersecurity update process, and applicable compliance plan \\
\midrule
Total programme: basis for quotation & Four gated stages with approximately 15 percent contingency & INR 10.5 to 19.7 lakh & Progression only after the previous evidence gate is met \\
\bottomrule
\end{tabularx}
\end{table*}

The first funded priority is the controlled validation rig rather than cosmetic product work. It resolves the largest scientific and commercial uncertainty: whether the alert remains specific after mounting, restart, load, speed, sensor, and motor identity change. Firmware and printed circuit board work can proceed in parallel only far enough to verify acquisition and replay parity. A field pilot should begin on noncritical assets and should keep every alert, inspection outcome, sensor intervention, and operating hour. This record is needed to set an actionable alert policy and to estimate service cost without extrapolating from balanced simulation cases.

The minimum product software consists of a versioned motor profile, acquisition self checks, feature extraction, a transparent baseline, an optional learned score, commissioning records, a local event queue, secure forwarding, and an audit trail. The dashboard should display the raw evidence behind an alert, the applicable operating cell, and the limits of the model. Fault names should be withheld when the measurements support only anomaly detection. This software scope is more valuable than adding unrelated features because it directly addresses the failure modes observed in the experiments.

\section{Conclusion}
SparkSense was reformulated as a motor profile guided monitoring framework with explicit evidence boundaries and a defined route to application. In controlled simulation, integer quantization preserved most decisions but did not protect amplitude based detection from gain and noise shifts. A load conditioned harmonic comparator cancelled common gain and detected injected drag and imbalance, but it was designed after inspection and remains unvalidated on measured multimodal motor data. Independent Ottawa vibration replay showed that the quantized model and a simple RMS threshold each detected 8 of 13 fully faulty recordings while generating 2 false alarms among 20 healthy recordings. The intended application is reviewable condition alerting for noncritical mixed motor fleets, especially where connectivity is intermittent. The next decisive experiment is synchronized acquisition of independently identified motor state, sensor state, current, speed, load, and vibration. The proposed four stage programme places that validation before field and preproduction claims. Until the evidence gates are closed, SparkSense should issue condition alerts without claiming fault classification, remaining life, or predictive maintenance outcomes.

\balance

\section*{Data Availability}
The Ottawa recordings are available through Mendeley Data, version 5, DOI: 10.17632/y2px5tg92h.5 \cite{sehri2023dataset}. The supplementary research package contains the simulation and replay code, parameters, virtual unit assignments, recording level decisions, and plotting scripts. Raw Ottawa recordings are excluded. No new physical measurements were collected for this study.

\section*{Declaration of AI Assistance}
AI tools assisted with code development, experiment execution, figure preparation, and manuscript drafting. The authors are responsible for verifying every analysis, statement, citation, and declaration before submission.

\begin{thebibliography}{14}
\bibitem{wheat2024} L. Wheat, M. von Mohrenschildt, S. R. Habibi, and D. Al Ani, ``Impact of data leakage in vibration signals used for bearing fault diagnosis,'' \emph{IEEE Access}, vol. 12, pp. 169879 to 169895, 2024, doi: 10.1109/ACCESS.2024.3497716.
\bibitem{jacob2018} B. Jacob \emph{et al.}, ``Quantization and training of neural networks for efficient integer arithmetic only inference,'' in \emph{Proc. IEEE Conf. Computer Vision and Pattern Recognition}, 2018, pp. 2704 to 2713, doi: 10.1109/CVPR.2018.00286.
\bibitem{arciniegas2025} S. Arciniegas \emph{et al.}, ``IoT device for detecting abnormal vibrations in motors using TinyML,'' \emph{Discover Internet of Things}, vol. 5, art. 41, 2025, doi: 10.1007/s43926-025-00142-4.
\bibitem{lei2020} Y. Lei, B. Yang, X. Jiang, F. Jia, N. Li, and A. K. Nandi, ``Applications of machine learning to machine fault diagnosis: A review and roadmap,'' \emph{Mechanical Systems and Signal Processing}, vol. 138, art. 106587, 2020, doi: 10.1016/j.ymssp.2019.106587.
\bibitem{iso20816} International Organization for Standardization, ``Mechanical vibration: Measurement and evaluation of machine vibration, Part 1: General guidelines,'' ISO 20816 1:2016, Geneva, Switzerland, 2016.
\bibitem{warden2019} P. Warden and D. Situnayake, \emph{TinyML: Machine Learning with TensorFlow Lite on Arduino and Ultra Low Power Microcontrollers}. Sebastopol, CA, USA: O'Reilly Media, 2019.
\bibitem{breiman2001} L. Breiman, ``Random forests,'' \emph{Machine Learning}, vol. 45, no. 1, pp. 5 to 32, 2001, doi: 10.1023/A:1010933404324.
\bibitem{adxl355} Analog Devices, ``ADXL354 and ADXL355 data sheet,'' rev. D, 2021. [Online]. Available: \url{https://www.analog.com/media/en/technical-documentation/data-sheets/adxl354_adxl355.pdf}.
\bibitem{sehri2023data} M. Sehri, P. Dumond, and M. Bouchard, ``University of Ottawa constant load and speed rolling element bearing vibration and acoustic fault signature datasets,'' \emph{Data in Brief}, vol. 49, art. 109327, 2023, doi: 10.1016/j.dib.2023.109327.
\bibitem{sehri2023dataset} M. Sehri and P. Dumond, ``University of Ottawa UORED VAFCLS bearing dataset,'' version 5, Mendeley Data, 2023, doi: 10.17632/y2px5tg92h.5.
\bibitem{gangsar2020} P. Gangsar and R. Tiwari, ``Signal based condition monitoring techniques for fault detection and diagnosis of induction motors: A state of the art review,'' \emph{Mechanical Systems and Signal Processing}, vol. 144, art. 106908, 2020, doi: 10.1016/j.ymssp.2020.106908.
\bibitem{zhao2019} R. Zhao, R. Yan, Z. Chen, K. Mao, P. Wang, and R. X. Gao, ``Deep learning and its applications to machine health monitoring,'' \emph{Mechanical Systems and Signal Processing}, vol. 115, pp. 213 to 237, 2019, doi: 10.1016/j.ymssp.2018.05.050.
\bibitem{david2021} R. David \emph{et al.}, ``TensorFlow Lite Micro: Embedded machine learning for TinyML systems,'' in \emph{Proceedings of Machine Learning and Systems}, vol. 3, pp. 800 to 811, 2021.
\bibitem{harris1978} F. J. Harris, ``On the use of windows for harmonic analysis with the discrete Fourier transform,'' \emph{Proceedings of the IEEE}, vol. 66, no. 1, pp. 51 to 83, 1978, doi: 10.1109/PROC.1978.10837.
\end{thebibliography}

\end{document}
